\emph{See also:} \url{experimental_setup.md} for the canonical hardware
and firmware stack (this chapter's ``Test conditions'' section
reproduces the subset relevant to the memcpy result),
\url{evaluation.md} for the cross-cutting results summary where this
chapter's findings appear as RQ2, \url{artifact.md} for the claim-to-CSV
map, and \url{statistical_notes.md} for the statistical conventions
applied to the per-run tables below.

\section{Abstract}\label{abstract}

The end-to-end vision pipeline on the SentAI board (Coral Dev Board
Micro, NXP i.MX RT1176) was stuck at \textasciitilde13 FPS even though
the camera delivers 15 FPS and the EdgeTPU invoke on our lightest model
only costs \textasciitilde50 ms. Instrumenting the pipeline's InferTask
exposed a single dominant cost: a 786 KB
\texttt{memcpy(tensor\_buf,\ staging\_buf,\ total)} from SDRAM to SDRAM
that consumed \textbf{24 ms per frame} due to CPU-driven cache-line
traffic through the SEMC controller. Replacing that one \texttt{memcpy}
with an eDMA memory-to-memory transfer configured for 32-byte AXI bursts
took the copy down to \textbf{14.6 ms} and raised sustained pipeline
throughput from \textbf{13.41 FPS} to \textbf{15.47 FPS} --- enough to
match the sensor rate.

\section{Test conditions}\label{test-conditions}

\subsection{Hardware}\label{hardware}

{\def\LTcaptype{none} % do not increment counter
\begin{longtable}[]{@{}
  >{\raggedright\arraybackslash}p{(\linewidth - 2\tabcolsep) * \real{0.6111}}
  >{\raggedright\arraybackslash}p{(\linewidth - 2\tabcolsep) * \real{0.3889}}@{}}
\toprule\noalign{}
\begin{minipage}[b]{\linewidth}\raggedright
Component
\end{minipage} & \begin{minipage}[b]{\linewidth}\raggedright
Value
\end{minipage} \\
\midrule\noalign{}
\endhead
\bottomrule\noalign{}
\endlastfoot
MCU & NXP i.MX RT1176 (Cortex-M7 @ 800 MHz + Cortex-M4) \\
On-chip EdgeTPU & Coral/Google TPU, internal USB2 bus \\
External RAM & 16 MB SDR-SDRAM on SEMC (166 MHz) \\
OCRAM / DTCM / ITCM & 1.25 MB / 256 KB / 256 KB \\
Camera & OV5640-based coralmicro module, 1280×720 native, 15 FPS
streaming \\
USB to host & CDC-ACM (REPL) + CDC-NCM (IP 10.0.0.1) \\
Board & Coral Dev Board Micro dev kit, powered via USB-C \\
\end{longtable}
}

\subsection{Firmware}\label{firmware}

{\def\LTcaptype{none} % do not increment counter
\begin{longtable}[]{@{}
  >{\raggedright\arraybackslash}p{(\linewidth - 2\tabcolsep) * \real{0.6111}}
  >{\raggedright\arraybackslash}p{(\linewidth - 2\tabcolsep) * \real{0.3889}}@{}}
\toprule\noalign{}
\begin{minipage}[b]{\linewidth}\raggedright
Parameter
\end{minipage} & \begin{minipage}[b]{\linewidth}\raggedright
Value
\end{minipage} \\
\midrule\noalign{}
\endhead
\bottomrule\noalign{}
\endlastfoot
Build & \texttt{sentai\_runtime} build \#622+ (linker script
\texttt{MIMXRT1176xxxxx\_cm7\_ram\_mp.ld}) \\
RTOS & FreeRTOS (CMSIS M7 build, 1 ms tick) \\
MicroPython GC heap & 512 KB, in \texttt{.sdram\_bss} \\
Pipeline tasks & \texttt{det\_prep} prio 2 · \texttt{det\_infer} prio
3 \\
Camera task & \texttt{camera\_task} prio
\texttt{configMAX\_PRIORITIES\ -\ 1} \\
\end{longtable}
}

\subsection{Model under test}\label{model-under-test}

{\def\LTcaptype{none} % do not increment counter
\begin{longtable}[]{@{}
  >{\raggedright\arraybackslash}p{(\linewidth - 2\tabcolsep) * \real{0.6111}}
  >{\raggedright\arraybackslash}p{(\linewidth - 2\tabcolsep) * \real{0.3889}}@{}}
\toprule\noalign{}
\begin{minipage}[b]{\linewidth}\raggedright
Parameter
\end{minipage} & \begin{minipage}[b]{\linewidth}\raggedright
Value
\end{minipage} \\
\midrule\noalign{}
\endhead
\bottomrule\noalign{}
\endlastfoot
File &
\texttt{/yolo\_1\_class\_512\_1\_upsample\_512\_inloc\_de\_1024\_la\_P5\_32.tflite} \\
On-flash size & 5 591 680 B (5.33 MB) \\
Input & \texttt{uint8{[}1,\ 512,\ 512,\ 3{]}} --- \textbf{786 432 B}
(this is the buffer we memcpy) \\
Input quantisation & scale = 1/255, zero\_point = 0 \\
Output & \texttt{uint8{[}1,\ 1344,\ 6{]}} --- 8 064 B ---
YOLOv5-enhanced anchor format \\
Architecture & YOLOv5 enhanced (single upsample at P5/32, 1-class
head) \\
Output quantisation & scale = 1/255, zero\_point = 0 \\
Output row semantics &
\texttt{{[}cx,\ cy,\ w,\ h,\ obj\_conf,\ class\_conf{]}} (normalised
0..1) \\
TFLite arena & 799 KB used / 8 192 KB available \\
EdgeTPU opened in & mode 3 (\texttt{kMax}) \\
\end{longtable}
}

\subsection{Scene}\label{scene}

\begin{itemize}
\tightlist
\item
  Static scene (camera does not move during measurement).
\item
  Before/after scene snapshots are saved by every experiment that
  touches the camera --- see
  \texttt{diag.snapshot\_scene("before"\textbar{}"after")} in
  \href{../diag/_session.py}{diag/\_session.py}. Stored at
  \texttt{/diags/\textless{}session\textgreater{}/scene\_cam\textless{}N\textgreater{}\_(before\textbar{}after)\_\textless{}W\textgreater{}x\textless{}H\textgreater{}.jpg}.
\item
  For the runs reported here the scene contained no target class
  instances, so \texttt{num\_detections\ ==\ 0} every frame --- the NMS
  execution path is still run in full (candidate scan over all 1344
  anchors), only the sort and IoU stages exit early.
\end{itemize}

\subsection{Methodology}\label{methodology}

\begin{itemize}
\tightlist
\item
  Each run = 20 measurement frames after 1 warm-up frame.
\item
  10 back-to-back runs per configuration to capture variance.
\item
  \textbf{Both configurations measured from the same firmware image}
  using the runtime A/B flag
  \texttt{sentai.pipeline.dma\_memcpy(0\textbar{}1)}. Switching between
  CPU memcpy and eDMA is one volatile store; no reflash, no reboot, no
  camera restart between the two blocks. This removes every confound
  that could come from firmware drift, thermal state, or camera
  calibration variance.
\item
  Timing collected via \texttt{sentai.pipeline.get\_ex()} which returns
  the firmware-measured \texttt{inference\_ms} (pure TPU
  \texttt{Invoke}), \texttt{memcpy\_ms} (staging → tensor copy),
  \texttt{nms\_ms} (post-processing), and \texttt{total\_ms} (full
  InferTask iteration). All reported in milliseconds.
\item
  Host (Python) wall-clock interval is the delta between consecutive
  \texttt{pipeline.get\_ex()} returns; the gap between
  \texttt{wall\_interval} and \texttt{infer\_total} measures everything
  outside InferTask (queue IPC, MicroPython overhead, mp\_repl wake-up).
\item
  \texttt{verbose(0)} for the whole loop body --- no \texttt{printf}
  traffic on CDC-ACM, so USB never saturates and host reads don't stall.
\item
  Benchmark driver: \href{../_e15_ab.py}{\_e15\_ab.py}.
\end{itemize}

\section{\texorpdfstring{Baseline --- CPU \texttt{memcpy} (SDRAM → cache
→
SDRAM)}{Baseline --- CPU memcpy (SDRAM → cache → SDRAM)}}\label{baseline-cpu-memcpy-sdram-cache-sdram}

Flag set by \texttt{sentai.pipeline.dma\_memcpy(0)}. The InferTask uses
a plain \texttt{memcpy()} to move the prepared frame into the TFLite
input tensor. This matches the state right after the camera-drain fix
(non-blocking \texttt{TryGetRawFrame}, documented separately in
\url{usb.md} and \url{camera.md}).

\subsection{Per-run results (10 × 20 frames, same firmware as
optimised)}\label{per-run-results-10-20-frames-same-firmware-as-optimised}

{\def\LTcaptype{none} % do not increment counter
\begin{longtable}[]{@{}
  >{\raggedleft\arraybackslash}p{(\linewidth - 12\tabcolsep) * \real{0.1136}}
  >{\raggedleft\arraybackslash}p{(\linewidth - 12\tabcolsep) * \real{0.1136}}
  >{\raggedleft\arraybackslash}p{(\linewidth - 12\tabcolsep) * \real{0.1364}}
  >{\raggedleft\arraybackslash}p{(\linewidth - 12\tabcolsep) * \real{0.1364}}
  >{\raggedleft\arraybackslash}p{(\linewidth - 12\tabcolsep) * \real{0.1136}}
  >{\raggedleft\arraybackslash}p{(\linewidth - 12\tabcolsep) * \real{0.2727}}
  >{\raggedleft\arraybackslash}p{(\linewidth - 12\tabcolsep) * \real{0.1136}}@{}}
\toprule\noalign{}
\begin{minipage}[b]{\linewidth}\raggedleft
Run
\end{minipage} & \begin{minipage}[b]{\linewidth}\raggedleft
wall
\end{minipage} & \begin{minipage}[b]{\linewidth}\raggedleft
invoke
\end{minipage} & \begin{minipage}[b]{\linewidth}\raggedleft
memcpy
\end{minipage} & \begin{minipage}[b]{\linewidth}\raggedleft
nms
\end{minipage} & \begin{minipage}[b]{\linewidth}\raggedleft
infer\_total
\end{minipage} & \begin{minipage}[b]{\linewidth}\raggedleft
FPS
\end{minipage} \\
\midrule\noalign{}
\endhead
\bottomrule\noalign{}
\endlastfoot
1 & 75.9 & 50.8 & 24.0 & 0.3 & 75.1 & 13.2 \\
2 & 74.8 & 49.9 & 24.4 & 0.4 & 74.6 & 13.4 \\
3 & 74.6 & 50.1 & 24.3 & 0.3 & 74.6 & 13.4 \\
4 & 74.7 & 50.3 & 24.0 & 0.4 & 74.6 & 13.4 \\
5 & 74.6 & 50.0 & 24.3 & 0.2 & 74.4 & 13.4 \\
6 & 74.2 & 50.0 & 24.1 & 0.1 & 74.3 & 13.5 \\
7 & 74.6 & 50.3 & 24.1 & 0.2 & 74.6 & 13.4 \\
8 & 74.9 & 50.3 & 24.1 & 0.2 & 74.7 & 13.4 \\
9 & 74.2 & 50.1 & 24.0 & 0.1 & 74.2 & 13.5 \\
10 & 74.4 & 50.3 & 24.0 & 0.2 & 74.4 & 13.4 \\
\textbf{mean} & \textbf{74.7} & \textbf{50.2} & \textbf{24.1} &
\textbf{0.2} & \textbf{74.5} & \textbf{13.39} \\
\end{longtable}
}

\textbf{FPS stats: mean 13.39, min 13.18, max 13.49, σ ≈ 0.10.}

\subsection{Interpretation}\label{interpretation}

The pipeline loop is fully dominated by InferTask:
\texttt{gap\ =\ wall\ −\ total\ ≈\ 0.2\ ms}, which is the cost of the
FreeRTOS \texttt{xQueueSend}, CDC-NCM wake-up and MicroPython tuple
allocation. Everything else is inside the firmware, and of that
\texttt{memcpy} alone is \textbf{32 \% of the critical path} (24.1 /
74.4).

Throughput during the copy = 786 432 B / 24 ms ≈ \textbf{32.8 MB/s}.
That is far below SDRAM's theoretical 400 MB/s peak. The CPU is
single-word- copying through the D-cache, which forces:

\begin{enumerate}
\def\labelenumi{\arabic{enumi}.}
\tightlist
\item
  Each 32-byte cache line of \textbf{source} is fetched from SDRAM (24
  576 fills for 786 KB).
\item
  Each 32-byte cache line of \textbf{destination} is \emph{also} fetched
  from SDRAM, because the Cortex-M7 D-cache is write-allocate.
\item
  The cache line is modified in place.
\item
  The dirty destination line is later written back to SDRAM.
\end{enumerate}

So the SEMC bus sees \textasciitilde2.4× the payload size
(\textasciitilde1.9 MB) and all of it as isolated single-beat AXI
accesses --- the bus cannot coalesce consecutive CPU stores into
back-to-back bursts. Application note AN12437 explicitly calls this out:
\emph{``SDRAM can reach high throughput when accessed by LCD and PXP, as
these two masters support back-to-back access, with better performance
compared to other master access, but dropping when accessed by CPU
core.''}

\section{Optimised --- eDMA memcpy with 32-byte AXI
bursts}\label{optimised-edma-memcpy-with-32-byte-axi-bursts}

The replacement helper lives in
\href{../detection_task.cc}{detection\_task.cc}:

\begin{verbatim}
// DMA0 channel 31 (audio driver uses ch 0).  One-time init.
static constexpr uint32_t kSentaiDmaChannel = 31;
static edma_handle_t s_dma_memcpy_handle;

// Pick the widest transfer width the addresses and size allow.
// 32-byte → one 8-beat AXI burst per request, exactly the pattern SEMC
// is optimised for.
uint32_t width = 4;
if ((((uintptr_t)src | (uintptr_t)dst | size) & 0x1Fu) == 0u)      width = 32;
else if ((((uintptr_t)src | (uintptr_t)dst | size) & 0x07u) == 0u) width = 8;

edma_transfer_config_t tcfg;
EDMA_PrepareTransfer(&tcfg, src, width, dst, width,
                     /*bytesEachRequest=*/size,
                     /*transferBytes=*/size,
                     kEDMA_MemoryToMemory);
EDMA_SubmitTransfer(&s_dma_memcpy_handle, &tcfg);
EDMA_StartTransfer(&s_dma_memcpy_handle);          // arm (SERQ)
EDMA_TriggerChannelStart(DMA0, kSentaiDmaChannel); // fire first minor loop

// Bounded polled wait — no ISR, no semaphore, no scheduler coupling.
const TickType_t deadline = xTaskGetTickCount() + pdMS_TO_TICKS(50);
while ((EDMA_GetChannelStatusFlags(DMA0, kSentaiDmaChannel)
        & kEDMA_DoneFlag) == 0) {
    if (xTaskGetTickCount() > deadline) return false;  // fallback to CPU memcpy
}
EDMA_ClearChannelStatusFlags(DMA0, kSentaiDmaChannel,
                             kEDMA_InterruptFlag | kEDMA_DoneFlag);
DCACHE_InvalidateByRange((uint32_t)dst, size);
\end{verbatim}

Two non-obvious details that cost measurable debug time:

\begin{enumerate}
\def\labelenumi{\arabic{enumi}.}
\tightlist
\item
  \textbf{\texttt{EDMA\_StartTransfer} only sets SERQ} (the
  request-enable bit). Memory-to-memory transfers have no peripheral
  request line, so the minor loop never fires until we explicitly
  software-trigger via \texttt{EDMA\_TriggerChannelStart} (SSRT
  register). Without that call the channel sits armed forever and the 50
  ms bounded wait expires.
\item
  \textbf{Poll \texttt{kEDMA\_DoneFlag}, not
  \texttt{kEDMA\_InterruptFlag}}. Interrupts are not enabled on this
  channel (no IRQ handler installed), so the interrupt flag would never
  set.
\end{enumerate}

No ISR is installed; the wait is a bounded polled loop. A 786 KB
transfer completes in \textasciitilde15 ms of wall time with a 35 ms
safety margin, so the CPU can't livelock. On polled-loop timeout (DMA
hardware stuck or mis-configured) the code falls back to plain
\texttt{memcpy} --- the pipeline degrades, doesn't brick.

\subsection{Cache strategy}\label{cache-strategy}

Both buffers are written exclusively by DMA masters in steady state
(\texttt{s\_staging\_buf} by PXP after camera capture; the TFLite input
tensor by our new eDMA). DMA writes bypass the D-cache, so neither
buffer holds dirty cache lines under normal operation. Our first attempt
called \texttt{DCACHE\_CleanInvalidateByRange} on both source and
destination before the DMA, for safety. That added
\textbf{\textasciitilde20 ms} per frame --- cache ops walk all address
ranges and on SDRAM the walk itself serialises with SEMC bus writebacks
of whatever cache lines happen to be dirty. We measured an end-to-end
memcpy of 78 ms with the defensive cache flush, worse than the original
CPU \texttt{memcpy}.

Removing the pre-DMA \texttt{CleanInvalidate} --- while keeping a
mandatory \texttt{DCACHE\_InvalidateByRange(dst,\ size)} \textbf{after}
the DMA so subsequent CPU reads see fresh bytes --- was what actually
delivered the win.

\subsection{Per-run results (10 × 20 frames, same firmware,
dma\_memcpy=1)}\label{per-run-results-10-20-frames-same-firmware-dma_memcpy1}

Flag set by \texttt{sentai.pipeline.dma\_memcpy(1)}. Measured
immediately after the baseline block --- same camera, same scene, same
model, same Python loop body, same TPU state. Only the runtime flag
changed.

{\def\LTcaptype{none} % do not increment counter
\begin{longtable}[]{@{}
  >{\raggedleft\arraybackslash}p{(\linewidth - 12\tabcolsep) * \real{0.1136}}
  >{\raggedleft\arraybackslash}p{(\linewidth - 12\tabcolsep) * \real{0.1136}}
  >{\raggedleft\arraybackslash}p{(\linewidth - 12\tabcolsep) * \real{0.1364}}
  >{\raggedleft\arraybackslash}p{(\linewidth - 12\tabcolsep) * \real{0.1364}}
  >{\raggedleft\arraybackslash}p{(\linewidth - 12\tabcolsep) * \real{0.1136}}
  >{\raggedleft\arraybackslash}p{(\linewidth - 12\tabcolsep) * \real{0.2727}}
  >{\raggedleft\arraybackslash}p{(\linewidth - 12\tabcolsep) * \real{0.1136}}@{}}
\toprule\noalign{}
\begin{minipage}[b]{\linewidth}\raggedleft
Run
\end{minipage} & \begin{minipage}[b]{\linewidth}\raggedleft
wall
\end{minipage} & \begin{minipage}[b]{\linewidth}\raggedleft
invoke
\end{minipage} & \begin{minipage}[b]{\linewidth}\raggedleft
memcpy
\end{minipage} & \begin{minipage}[b]{\linewidth}\raggedleft
nms
\end{minipage} & \begin{minipage}[b]{\linewidth}\raggedleft
infer\_total
\end{minipage} & \begin{minipage}[b]{\linewidth}\raggedleft
FPS
\end{minipage} \\
\midrule\noalign{}
\endhead
\bottomrule\noalign{}
\endlastfoot
1 & 64.3 & 49.5 & 14.6 & 0.2 & 64.3 & 15.6 \\
2 & 66.3 & 50.2 & 14.6 & 0.3 & 65.1 & 15.1 \\
3 & 64.6 & 49.7 & 14.6 & 0.2 & 64.4 & 15.5 \\
4 & 65.9 & 49.8 & 14.6 & 0.2 & 64.7 & 15.2 \\
5 & 64.3 & 49.5 & 14.6 & 0.2 & 64.2 & 15.6 \\
6 & 66.1 & 50.1 & 14.4 & 0.2 & 64.7 & 15.1 \\
7 & 64.7 & 49.7 & 14.6 & 0.2 & 64.5 & 15.5 \\
8 & 66.1 & 50.0 & 14.6 & 0.3 & 64.8 & 15.1 \\
9 & 64.6 & 49.7 & 14.5 & 0.2 & 64.4 & 15.5 \\
10 & 66.1 & 49.9 & 14.7 & 0.2 & 64.8 & 15.1 \\
\textbf{mean} & \textbf{65.3} & \textbf{49.8} & \textbf{14.6} &
\textbf{0.2} & \textbf{64.6} & \textbf{15.32} \\
\end{longtable}
}

\textbf{FPS stats: mean 15.32, min 15.07, max 15.56, σ ≈ 0.22.}

Throughput during the copy = 786 432 B / 14.6 ms ≈ \textbf{53.9 MB/s}
--- a \textbf{1.65× improvement} for the same buffer size and addresses.

The observed pipeline FPS sits right at the camera sensor rate,
confirming that with the DMA optimisation engaged the sensor --- not the
compute path --- is the rate-limiting step.

\section{Comparison and discussion}\label{comparison-and-discussion}

\subsection{Per-stage cost (paired A/B on same firmware, 200 frames
each)}\label{per-stage-cost-paired-ab-on-same-firmware-200-frames-each}

{\def\LTcaptype{none} % do not increment counter
\begin{longtable}[]{@{}
  >{\raggedright\arraybackslash}p{(\linewidth - 6\tabcolsep) * \real{0.2000}}
  >{\raggedleft\arraybackslash}p{(\linewidth - 6\tabcolsep) * \real{0.2667}}
  >{\raggedleft\arraybackslash}p{(\linewidth - 6\tabcolsep) * \real{0.2667}}
  >{\raggedleft\arraybackslash}p{(\linewidth - 6\tabcolsep) * \real{0.2667}}@{}}
\toprule\noalign{}
\begin{minipage}[b]{\linewidth}\raggedright
Stage
\end{minipage} & \begin{minipage}[b]{\linewidth}\raggedleft
CPU memcpy
\end{minipage} & \begin{minipage}[b]{\linewidth}\raggedleft
eDMA 32-B bursts
\end{minipage} & \begin{minipage}[b]{\linewidth}\raggedleft
Δ
\end{minipage} \\
\midrule\noalign{}
\endhead
\bottomrule\noalign{}
\endlastfoot
Invoke (TPU) & 50.2 ms & 49.8 ms & −0.4 ms (noise) \\
staging → tensor copy & \textbf{24.1 ms} & \textbf{14.6 ms} &
\textbf{−9.5 ms (−39 \%)} \\
NMS (1-class × 1344 anchors, YOLOv5 layout) & 0.2 ms & 0.2 ms & --- \\
InferTask total & 74.5 ms & 64.6 ms & \textbf{−9.9 ms} \\
wall interval & 74.7 ms & 65.3 ms & −9.4 ms \\
\end{longtable}
}

The saving is entirely explained by the copy; invoke, NMS and the
Python/IPC gap are unchanged within noise. Importantly, \texttt{invoke}
did \textbf{not} get slower even though the eDMA now competes with the
TPU's internal USB bus for SEMC bandwidth --- the transfer is short
enough (15 ms, before invoke starts) that it finishes before invoke's
input upload begins.

\subsection{Throughput headline}\label{throughput-headline}

{\def\LTcaptype{none} % do not increment counter
\begin{longtable}[]{@{}
  >{\raggedright\arraybackslash}p{(\linewidth - 8\tabcolsep) * \real{0.1579}}
  >{\raggedleft\arraybackslash}p{(\linewidth - 8\tabcolsep) * \real{0.2105}}
  >{\raggedleft\arraybackslash}p{(\linewidth - 8\tabcolsep) * \real{0.2105}}
  >{\raggedleft\arraybackslash}p{(\linewidth - 8\tabcolsep) * \real{0.2105}}
  >{\raggedleft\arraybackslash}p{(\linewidth - 8\tabcolsep) * \real{0.2105}}@{}}
\toprule\noalign{}
\begin{minipage}[b]{\linewidth}\raggedright
Config
\end{minipage} & \begin{minipage}[b]{\linewidth}\raggedleft
FPS mean
\end{minipage} & \begin{minipage}[b]{\linewidth}\raggedleft
FPS min
\end{minipage} & \begin{minipage}[b]{\linewidth}\raggedleft
FPS max
\end{minipage} & \begin{minipage}[b]{\linewidth}\raggedleft
Δ vs baseline
\end{minipage} \\
\midrule\noalign{}
\endhead
\bottomrule\noalign{}
\endlastfoot
CPU memcpy (\texttt{dma\_memcpy(0)}) & 13.39 & 13.18 & 13.49 & --- \\
eDMA 32-B bursts (\texttt{dma\_memcpy(1)}) & \textbf{15.32} &
\textbf{15.07} & \textbf{15.56} & \textbf{+14.4 \%} \\
Camera sensor rate (hard ceiling) & 15.00 & --- & --- & --- \\
\end{longtable}
}

The optimised configuration operates right at the camera sensor rate (15
FPS ≈ 66.7 ms/frame); individual measurements slightly exceed this when
the pipeline drains a small in-flight queue over one cycle and catches
up on the next. Sustained over a longer window the effective rate is
bounded by the camera. The 1.92 FPS absolute gain corresponds to a
\textbf{14.4 \% relative throughput improvement} measured under
controlled A/B conditions in the same firmware image.

\subsection{Variance}\label{variance}

σ grows from 0.09 to 0.20 FPS --- not material for the use case but
worth noting. The extra variance comes from the 50 ms DMA polling path
occasionally getting preempted by the camera task or the watchdog task;
this would disappear if we re-armed the polling using a
task-notification wake-up driven by the eDMA completion IRQ, at the cost
of additional ISR state. We judged the trade-off not worth it while we
are comfortably above the sensor rate.

\section{Why we chose eDMA over
alternatives}\label{why-we-chose-edma-over-alternatives}

{\def\LTcaptype{none} % do not increment counter
\begin{longtable}[]{@{}
  >{\raggedright\arraybackslash}p{(\linewidth - 4\tabcolsep) * \real{0.3000}}
  >{\raggedleft\arraybackslash}p{(\linewidth - 4\tabcolsep) * \real{0.4000}}
  >{\raggedright\arraybackslash}p{(\linewidth - 4\tabcolsep) * \real{0.3000}}@{}}
\toprule\noalign{}
\begin{minipage}[b]{\linewidth}\raggedright
Option
\end{minipage} & \begin{minipage}[b]{\linewidth}\raggedleft
Expected time
\end{minipage} & \begin{minipage}[b]{\linewidth}\raggedright
Why we didn't pick it
\end{minipage} \\
\midrule\noalign{}
\endhead
\bottomrule\noalign{}
\endlastfoot
Keep CPU \texttt{memcpy} + \texttt{\_\_builtin\_prefetch} hints &
\textasciitilde18 ms & Only \textasciitilde25 \% saving, still CPU-bound
on SEMC \\
Move tensor to OCRAM & N/A & TFLite arena is 8 MB, exceeds 1.25 MB
OCRAM \\
Move staging to OCRAM & mixed & Halves read-side cost but write-side
still on SEMC; net \textasciitilde10 ms, small gain for a large
refactor \\
PXP with PS→output 1:1 & \textasciitilde3-5 ms & PXP is already used by
PrepTask; serialising PXP across PrepTask+InferTask would defeat
pipeline parallelism \\
\textbf{eDMA mem-to-mem} & \textbf{\textasciitilde14-15 ms} &
\textbf{Independent engine, back-to-back SEMC bursts, no conflict with
PXP, isolated to a single helper function} \\
\end{longtable}
}

PXP would in principle be even faster than eDMA (AN12437 lists PXP as
best-in-class for SDRAM throughput), but PXP is the camera-side resize
engine and scheduling it twice per frame pulls PrepTask and InferTask
into a shared-resource dance that kills parallelism. eDMA channel 31 is
independent of every other user in the firmware, so the change stays
local and analysable.

\section{Failure modes}\label{failure-modes}

Per \href{../agent/embeded.md}{agent/embeded.md}: every optimisation
must have a defined failure path.

{\def\LTcaptype{none} % do not increment counter
\begin{longtable}[]{@{}
  >{\raggedright\arraybackslash}p{(\linewidth - 4\tabcolsep) * \real{0.3333}}
  >{\raggedright\arraybackslash}p{(\linewidth - 4\tabcolsep) * \real{0.3333}}
  >{\raggedright\arraybackslash}p{(\linewidth - 4\tabcolsep) * \real{0.3333}}@{}}
\toprule\noalign{}
\begin{minipage}[b]{\linewidth}\raggedright
Failure
\end{minipage} & \begin{minipage}[b]{\linewidth}\raggedright
Detection
\end{minipage} & \begin{minipage}[b]{\linewidth}\raggedright
Action
\end{minipage} \\
\midrule\noalign{}
\endhead
\bottomrule\noalign{}
\endlastfoot
Addresses not 4-byte aligned, or \texttt{size\ \%\ 4\ !=\ 0} &
compile-time alignment check fails & \texttt{dma\_ok\ =\ false} → falls
back to CPU \texttt{memcpy} \\
eDMA submit returns non-success & Return value of
\texttt{EDMA\_SubmitTransfer} & Return \texttt{false} → CPU
\texttt{memcpy} \\
Channel stuck (DONE flag never sets) & 50 ms bounded polled wait expires
& Return \texttt{false} → CPU \texttt{memcpy} \\
Some other task takes DMA0 ch31 & Not detected; documented single-owner
assumption & ch31 is dedicated to this helper --- audio uses ch0; camera
doesn't use eDMA \\
\end{longtable}
}

No silent corruption: the DMA writes to an empty buffer the TFLite
tensor alone will read; \texttt{DCACHE\_InvalidateByRange(dst)} after
the DMA guarantees the next invoker reads fresh bytes regardless of
cache state.

\section{Reproducing these results}\label{reproducing-these-results}

\begin{verbatim}
# On the device REPL, after flashing build #622+.
# One-shot E15 run — creates /diags/sNNN_e15_512/ with CSV + scene_{before,after}.jpg
import diag
diag.e15_pipeline_parallel_512(repetitions=5)

# 10x benchmark table (dma_memcpy stays at whatever the flag is set to)
import sentai
exec(sentai.fs.read_str('/_e15_table.py'))

# *** Paired A/B benchmark (this document's tables come from here) ***
# Runs 10x20 frames with DMA off, then 10x20 frames with DMA on,
# prints a side-by-side summary + FPS improvement.
exec(sentai.fs.read_str('/_e15_ab.py'))

# Manual toggle at any time:
sentai.pipeline.dma_memcpy(0)   # CPU memcpy (baseline)
sentai.pipeline.dma_memcpy(1)   # eDMA 32-B bursts (optimised)
sentai.pipeline.dma_memcpy()    # read current setting
\end{verbatim}

Every run writes:

\begin{itemize}
\tightlist
\item
  \texttt{/diags/sNNN\_e15\_512/manifest.csv} --- experiment log
\item
  \texttt{/diags/sNNN\_e15\_512/001\_e14\_pipeline\_par\_cam0\_512x512.csv}
  --- per-frame timing
\item
  \texttt{/diags/sNNN\_e15\_512/001\_e14\_pipeline\_par\_cam0\_512x512.txt}
  --- column key + params
\item
  \texttt{/diags/sNNN\_e15\_512/scene\_cam0\_before\_512x512.jpg} ---
  scene at experiment start
\item
  \texttt{/diags/sNNN\_e15\_512/scene\_cam0\_after\_512x512.jpg} ---
  scene at experiment end
\item
  \texttt{/diags/sNNN\_e15\_512/summary.txt} --- heap + duration summary
\end{itemize}

The pair of scene snapshots lets an operator confirm offline that the
camera was pointing at the expected target and that nothing moved during
the run. This matters when \texttt{num\_detections\ ==\ 0} --- it
distinguishes \emph{``scene had no target''} from \emph{``detector
missed''}.

\section{Cross-model comparison --- E14 (COCO 80-class) vs E15
(1-class)}\label{cross-model-comparison-e14-coco-80-class-vs-e15-1-class}

To put the memcpy optimisation in context we ran the pipeline 20 times
per experiment (20 frames each, 400 frames total per configuration),
each pair of runs writing CSVs + scene snapshots under its own
\texttt{/diags/sNNN\_...} session for later offline analysis.

{\def\LTcaptype{none} % do not increment counter
\begin{longtable}[]{@{}
  >{\raggedright\arraybackslash}p{(\linewidth - 4\tabcolsep) * \real{0.3333}}
  >{\raggedright\arraybackslash}p{(\linewidth - 4\tabcolsep) * \real{0.3333}}
  >{\raggedright\arraybackslash}p{(\linewidth - 4\tabcolsep) * \real{0.3333}}@{}}
\toprule\noalign{}
\begin{minipage}[b]{\linewidth}\raggedright
Property
\end{minipage} & \begin{minipage}[b]{\linewidth}\raggedright
E14 (Ultralytics YOLOv8 COCO)
\end{minipage} & \begin{minipage}[b]{\linewidth}\raggedright
E15 (custom YOLOv5-enhanced, 1-class)
\end{minipage} \\
\midrule\noalign{}
\endhead
\bottomrule\noalign{}
\endlastfoot
Model file & \texttt{/yolo26n.edgetpu\_1.tflite} &
\texttt{/yolo\_1\_class\_512\_1\_upsample\_...\_P5\_32.tflite} \\
File size & 4.41 MB & 5.33 MB \\
Architecture family & YOLOv8 nano (Ultralytics) & YOLOv5 enhanced
(single upsample at P5/32) \\
Output layout & \texttt{{[}1,\ 84,\ 2100{]}} (v8 transposed) &
\texttt{{[}1,\ 1344,\ 6{]}} (v5-style:
\texttt{{[}cx,cy,w,h,obj,cls{]}}) \\
Num. classes & 80 (COCO) & 1 \\
Input shape & int8 \texttt{{[}1,\ 320,\ 320,\ 3{]}} --- 307 200 B &
uint8 \texttt{{[}1,\ 512,\ 512,\ 3{]}} --- 786 432 B \\
Output shape & int8 \texttt{{[}1,\ 84,\ 2100{]}} --- 176 400 B & uint8
\texttt{{[}1,\ 1344,\ 6{]}} --- 8 064 B \\
NMS inner cost & \textbf{80 classes × 2100 anchors = 168 000
dequant+compares/frame} & 1 class × 1344 anchors = 1 344
compares/frame \\
DMA memcpy & on (small 307 KB payload, marginal win) & on (decisive 786
KB payload) \\
Session & \texttt{s031\_e14\_x20/} & \texttt{s032\_e15\_x20/} \\
\end{longtable}
}

\subsection{Aggregated results (20 runs × 20 frames each,
verbose=0)}\label{aggregated-results-20-runs-20-frames-each-verbose0}

{\def\LTcaptype{none} % do not increment counter
\begin{longtable}[]{@{}lrcrc@{}}
\toprule\noalign{}
Metric & E14 mean & E14 range & E15 mean & E15 range \\
\midrule\noalign{}
\endhead
\bottomrule\noalign{}
\endlastfoot
frame\_interval (ms) & 156.9 & 156.1--157.8 & 69.4 & 64.3--70.9 \\
pipeline FPS (wall) & 6.37 & 6.34--6.41 & 14.40 & 14.10--15.53 \\
firmware avg\_fps & 6.19 & 6.00--6.20 & 14.08 & 13.40--14.30 \\
detections (sum) & 2 & --- & 0 & --- \\
\end{longtable}
}

The \textbf{E14 experiment is NMS-bound} on this dataset: the 80-class
yolo iteration pays \textasciitilde168 000 dequant + argmax operations
per frame, which on Cortex-M7 FPU at 800 MHz is visibly slow compared to
the 1-class hot path. E15's NMS sees only 1 344 single-class comparisons
per frame and runs in \textless{} 1 ms. Together with the much heavier
786 KB tensor copy that E15 would have to do if we hadn't switched to
eDMA, this is why the bigger model can still run more than 2× faster
than the smaller one.

\subsection{Detection of yolo layout \& class
count}\label{detection-of-yolo-layout-class-count}

To automate experiment setup across models with different class budgets,
the firmware now exposes \texttt{sentai.tpu.yolo\_info()} which returns
\texttt{(layout\_str,\ num\_classes,\ num\_anchors)} inferred purely
from the loaded model's output tensor shape (the edgetpu compiler strips
most metadata buffers from the .tflite binary, so shape is the only
reliable signal). Verified on both experiments:

{\def\LTcaptype{none} % do not increment counter
\begin{longtable}[]{@{}ll@{}}
\toprule\noalign{}
Model file & \texttt{yolo\_info()} \\
\midrule\noalign{}
\endhead
\bottomrule\noalign{}
\endlastfoot
yolo26n.edgetpu\_1.tflite (E14) &
\texttt{(\textquotesingle{}v8\textquotesingle{},\ 80,\ 2100)} \\
yolo\_1\_class\_512\_1\_upsample\_\ldots\_P5\_32 (E15) &
\texttt{(\textquotesingle{}v5\_like\textquotesingle{},\ 1,\ 1344)} \\
\end{longtable}
}

Source: \href{../sentai_runtime.cc}{\texttt{yolo\_infer\_info()} in
sentai\_runtime.cc}.

\subsection{Artefacts saved for offline
inspection}\label{artefacts-saved-for-offline-inspection}

Each run in a session writes:

\begin{itemize}
\tightlist
\item
  \texttt{NNN\_\textless{}tag\textgreater{}\_pipeline\_par\_cam0\_\textless{}W\textgreater{}x\textless{}H\textgreater{}.csv}
  --- per-frame wall\_interval, prep\_stall, infer\_stall,
  num\_detections
\item
  \texttt{NNN\_\textless{}tag\textgreater{}\_pipeline\_par\_cam0\_\textless{}W\textgreater{}x\textless{}H\textgreater{}.txt}
  --- column key + parameters (conf, iou, repetitions, frames\_received,
  fw\_processed, fw\_dropped)
\item
  \texttt{scene\_cam0\_before\_\textless{}W\textgreater{}x\textless{}H\textgreater{}.jpg}
  --- scene the operator pointed at, at session start. Lets offline
  inspection confirm what the model saw.
\item
  \texttt{scene\_cam0\_after\_\textless{}W\textgreater{}x\textless{}H\textgreater{}.jpg}
  --- same scene, session end. Paired with BEFORE for drift / movement
  verification.
\item
  \texttt{manifest.csv} + \texttt{summary.txt} --- per-session roll-up.
\end{itemize}

\section{Progression across the pipeline optimisation
effort}\label{progression-across-the-pipeline-optimisation-effort}

{\def\LTcaptype{none} % do not increment counter
\begin{longtable}[]{@{}
  >{\raggedright\arraybackslash}p{(\linewidth - 8\tabcolsep) * \real{0.2857}}
  >{\raggedleft\arraybackslash}p{(\linewidth - 8\tabcolsep) * \real{0.2286}}
  >{\raggedleft\arraybackslash}p{(\linewidth - 8\tabcolsep) * \real{0.1714}}
  >{\raggedleft\arraybackslash}p{(\linewidth - 8\tabcolsep) * \real{0.1429}}
  >{\raggedright\arraybackslash}p{(\linewidth - 8\tabcolsep) * \real{0.1714}}@{}}
\toprule\noalign{}
\begin{minipage}[b]{\linewidth}\raggedright
Revision
\end{minipage} & \begin{minipage}[b]{\linewidth}\raggedleft
memcpy
\end{minipage} & \begin{minipage}[b]{\linewidth}\raggedleft
wall
\end{minipage} & \begin{minipage}[b]{\linewidth}\raggedleft
FPS
\end{minipage} & \begin{minipage}[b]{\linewidth}\raggedright
Note
\end{minipage} \\
\midrule\noalign{}
\endhead
\bottomrule\noalign{}
\endlastfoot
Baseline (camera drain polls 4 s) & 24 ms & \textasciitilde200 ms & 5.0
& \texttt{TryGetRawFrame} not truly non-blocking \\
Non-blocking camera drain & 24 ms & 75 ms & 13.4 & Fix \#1 --- see
usb.md / camera.md \\
\textbf{+ eDMA memcpy} & \textbf{14.6 ms} & \textbf{64.7 ms} &
\textbf{15.5} & \textbf{Fix \#2 --- this document} \\
\end{longtable}
}

End-to-end, \textbf{3.1× speedup over the initial firmware}, reaching
the camera sensor rate.

\section{What's next}\label{whats-next}

\begin{itemize}
\tightlist
\item
  \texttt{invoke} at 49.6 ms now dominates the loop. The remaining
  headroom is inside \texttt{Invoke()} itself --- 28 ms cold when
  measured outside the pipeline, 49 ms steady-state while PrepTask runs
  concurrently, suggesting SEMC/USB bus contention with the TPU's
  internal DMA. Pushing beyond 15 FPS would require either a lighter
  model (smaller input → less TPU bus), or double-buffering the TFLite
  tensor so NMS/Invoke overlap more aggressively.
\item
  The \texttt{sentai\_dma\_memcpy} helper is local to
  \texttt{detection\_task.cc}. Three other places in the firmware copy
  large SDRAM buffers (\texttt{sentai\_cam\_capture\_rgb},
  \texttt{sentai\_cam\_capture\_jpeg},
  \texttt{sentai\_cam\_to\_tensor\_ex}). Once lifted into
  \texttt{libs/base/dma\_memcpy.*} the same win applies there.
\item
  Switch to completion-interrupt + FreeRTOS notification once any
  concurrent user contests channel 31; the current polled loop is only
  safe because no other task uses this channel.
\end{itemize}

\section{Sources}\label{sources}

\begin{itemize}
\tightlist
\item
  NXP \emph{AN12437 --- i.MX RT Series Performance Optimization} (SDRAM
  access patterns per master; PXP/LCD/eDMA back-to-back bursts).
\item
  NXP \emph{i.MX RT1170 Reference Manual} (\texttt{IMXRT1170RM}), eDMA
  and SEMC chapters (TCD layout, SSRT/SERQ semantics, AXI burst
  configuration).
\item
  Local code: \href{../detection_task.cc}{detection\_task.cc},
  \href{../modsentai_pipeline.c}{modsentai\_pipeline.c},
  \href{../diag/e_pipeline.py}{diag/e\_pipeline.py},
  \href{../diag/_session.py}{diag/\_session.py}.
\end{itemize}
